\documentclass[a4paper,12pt]{article}
\usepackage[T1]{fontenc}
\usepackage{amsmath}
\usepackage{amssymb}
\usepackage[english,french]{babel}
\usepackage{color}
\usepackage{footmisc}
\usepackage{graphicx}
%\usepackage{mathpazo}
\usepackage{multicol}
\usepackage{stmaryrd}
%\usepackage[scaled=.85]{beramono}
\usepackage{../iman,../pdfsetup}

%\oddsidemargin=4.6mm
%\evensidemargin=4.6mm
%\textwidth=150mm
%\topmargin=4.6mm
%\headheight=0mm
%\headsep=0mm
%\textheight=234mm

\def\Colon{\mathord{:\mkern-1.5mu:}}
%\def\lbrakk{\mathopen{\lbrack\mkern-3.25mu\lbrack}}
%\def\rbrakk{\mathclose{\rbrack\mkern-3.255mu\rbrack}}
\def\lparr{\mathopen{(\mkern-4mu\mid}}
\def\rparr{\mathclose{\mid\mkern-4mu)}}

\def\unk{{?}}
\def\undef{(\lambda x.\; \unk)}
%\def\unr{\textit{others}}
\def\unr{\ldots}
\def\Abs#1{\hbox{\rm{\flqq}}{\,#1\,}\hbox{\rm{\frqq}}}
\def\Q{{\smash{\lower.2ex\hbox{$\scriptstyle?$}}}}

\urlstyle{tt}

\begin{document}

\selectlanguage{english}

\title{\includegraphics[scale=0.5]{isabelle_sledgehammer} \\[4ex]
Hammering Away \\[\smallskipamount]
\Large A User's Guide to Sledgehammer for Isabelle/HOL}
\author{\hbox{} \\
Jasmin Christian Blanchette \\
{\normalsize Institut f\"ur Informatik, Technische Universit\"at M\"unchen} \\
\hbox{}}

\maketitle

\tableofcontents

\setlength{\parskip}{.7em plus .2em minus .1em}
\setlength{\parindent}{0pt}
\setlength{\abovedisplayskip}{\parskip}
\setlength{\abovedisplayshortskip}{.9\parskip}
\setlength{\belowdisplayskip}{\parskip}
\setlength{\belowdisplayshortskip}{.9\parskip}

% General-purpose enum environment with correct spacing
\newenvironment{enum}%
    {\begin{list}{}{%
        \setlength{\topsep}{.1\parskip}%
        \setlength{\partopsep}{.1\parskip}%
        \setlength{\itemsep}{\parskip}%
        \advance\itemsep by-\parsep}}
    {\end{list}}

\def\pre{\begingroup\vskip0pt plus1ex\advance\leftskip by\leftmargin
\advance\rightskip by\leftmargin}
\def\post{\vskip0pt plus1ex\endgroup}

\def\prew{\pre\advance\rightskip by-\leftmargin}
\def\postw{\post}

\section{Introduction}
\label{introduction}

Sledgehammer is a tool that applies first-order automatic theorem provers (ATPs)
on the current goal. The supported ATPs are E \cite{schulz-2002}, SPASS
\cite{weidenbach-et-al-2009}, and Vampire \cite{riazanov-voronkov-2002}, which
can be run locally or remotely via the SystemOnTPTP web service
\cite{sutcliffe-2000}.

The problem passed to ATPs consists of your current goal together with a
heuristic selection of hundreds of facts (theorems) from the current theory
context, filtered by relevance. Because jobs are run in the background, you can
continue to work on your proof by other means. Provers can be run in parallel.
Any reply (which may arrive minutes later) will appear in the Proof General
response buffer.

The result of a successful ATP proof search is some source text that usually
(but not always) reconstructs the proof within Isabelle, without requiring the
ATPs again. The reconstructed proof relies on the general-purpose Metis prover
\cite{metis}, which is fully integrated into Isabelle/HOL, with explicit
inferences going through the kernel. Thus its results are correct by
construction.

\newbox\boxA
\setbox\boxA=\hbox{\texttt{nospam}}

Examples of Sledgehammer use can be found in Isabelle's
\texttt{src/HOL/Metis\_Examples} directory.
Comments and bug reports concerning Sledgehammer or this manual should be
directed to
\texttt{blan{\color{white}nospam}\kern-\wd\boxA{}chette@\allowbreak
in.\allowbreak tum.\allowbreak de}.

\vskip2.5\smallskipamount

%\textbf{Acknowledgment.} The author would like to thank Mark Summerfield for
%suggesting several textual improvements.

\section{Installation}
\label{installation}

Sledgehammer is part of Isabelle, so you don't need to install it. However, it
relies on third-party automatic theorem provers (ATPs). Currently, E, SPASS, and
Vampire are supported. All of these are available remotely via SystemOnTPTP
\cite{sutcliffe-2000}, but if you want better performance you will need to
install at least E and SPASS locally.

There are three main ways to install ATPs on your machine:

\begin{enum}
\item[$\bullet$] If you installed an official Isabelle package with everything
inside, it should already include properly setup executables for E and SPASS,
ready to use.%
\footnote{Vampire's license prevents us from doing the same for this otherwise
wonderful tool.}

\item[$\bullet$] Alternatively, you can download the Isabelle-aware E and SPASS
binary packages from Isabelle's download page. Extract the archives, then add a
line to your \texttt{\char`\~/.isabelle/etc/components} file with the absolute path to
E or SPASS. For example, if the \texttt{components} does not exist
yet and you extracted SPASS to \texttt{/usr/local/spass-3.7}, create
the \texttt{components} file with the single line

\prew
\texttt{/usr/local/spass-3.7}
\postw

in it.

\item[$\bullet$] If you prefer to build E or SPASS yourself, or obtained a
Vampire executable from somewhere (e.g., \url{http://www.vprover.org/}),
set the environment variable \texttt{E\_HOME}, \texttt{SPASS\_HOME}, or
\texttt{VAMPIRE\_HOME} to the directory that contains the \texttt{eproof},
\texttt{SPASS}, or \texttt{vampire} executable. Sledgehammer has been tested
with E 1.0 and 1.2, SPASS 3.5 and 3.7, and Vampire 1.0%
\footnote{Following the rewrite of Vampire, the counter for version numbers was
reset to 0; hence the new Vampire 1.0 is more recent than Vampire 11.5.}%
. Since the ATPs' output formats are neither documented nor stable, other
versions of the ATPs might or might not work well with Sledgehammer.
\end{enum}

To check whether E and SPASS are installed, follow the example in
\S\ref{first-steps}.

Remote ATP invocation via the SystemOnTPTP web service requires Perl with the
World Wide Web Library (\texttt{libwww-perl}) installed.

\section{First Steps}
\label{first-steps}

To illustrate Sledgehammer in context, let us start a theory file and
attempt to prove a simple lemma:

\prew
\textbf{theory}~\textit{Scratch} \\
\textbf{imports}~\textit{Main} \\
\textbf{begin} \\[2\smallskipamount]
%
\textbf{lemma} ``$[a] = [b] \,\longleftrightarrow\, a = b$'' \\
\textbf{sledgehammer}
\postw

Instead of issuing the \textbf{sledgehammer} command, you can also find
Sledgehammer in the ``Commands'' submenu of the ``Isabelle'' menu in Proof
General or press the Emacs key sequence C-c C-a C-s.
Either way, Sledgehammer produces the following output after a few seconds:

\prew
\slshape
Sledgehammer: ATP ``\textit{e}'' for subgoal 1: \\
$([a] = [b]) = (a = b)$ \\
Try this command: \textbf{by} (\textit{metis hd.simps}). \\
To minimize the number of lemmas, try this: \\
\textbf{sledgehammer} \textit{minimize} [\textit{atp} = \textit{e}] (\textit{hd.simps}). \\[3\smallskipamount]
%
Sledgehammer: ATP ``\textit{spass}'' for subgoal 1: \\
$([a] = [b]) = (a = b)$ \\
Try this command: \textbf{by} (\textit{metis insert\_Nil last\_ConsL}). \\
To minimize the number of lemmas, try this: \\
\textbf{sledgehammer} \textit{minimize} [\textit{atp} = \textit{spass}] (\textit{insert\_Nil last\_ConsL}). \\[3\smallskipamount]
%
Sledgehammer: ATP ``\textit{remote\_vampire}'' for subgoal 1: \\
$([a] = [b]) = (a = b)$ \\
Try this command: \textbf{by} (\textit{metis One\_nat\_def\_raw empty\_replicate} \\
\phantom{Try this command: \textbf{by} (\textit{metis~}}\textit{insert\_Nil last\_ConsL replicate\_Suc}). \\
To minimize the number of lemmas, try this: \\
\textbf{sledgehammer} \textit{minimize} [\textit{atp} = \textit{remote\_vampire}] \\
\phantom{\textbf{sledgehammer}~}(\textit{One\_nat\_def\_raw empty\_replicate insert\_Nil} \\
\phantom{\textbf{sledgehammer}~(}\textit{last\_ConsL replicate\_Suc}).
\postw

Sledgehammer ran E, SPASS, and the remote version of Vampire in parallel. If E
is not installed (\S\ref{installation}), you will see references to
its remote American cousin \textit{remote\_e} instead of
\textit{e}; and if SPASS is not installed, it will not appear in the output.

Based on each ATP proof, Sledgehammer gives a one-liner proof that uses the
\textit{metis} method. You can click them and insert them into the theory text.
You can click the ``\textbf{sledgehammer} \textit{minimize}'' command if you
want to look for a shorter (and faster) proof. But here the proof found by E
looks perfect, so click it to finish the proof.

You can ask Sledgehammer for an Isar text proof by passing the
\textit{isar\_proof} option:

\prew
\textbf{sledgehammer} [\textit{isar\_proof}]
\postw

When Isar proof construction is successful, it can yield proofs that are more
readable and also faster than the \textit{metis} one-liners. This feature is
experimental.

\section{Hints}
\label{hints}

For best results, first simplify your problem by calling \textit{auto} or at
least \textit{safe} followed by \textit{simp\_all}. None of the ATPs contain
arithmetic decision procedures. They are not especially good at heavy rewriting,
but because they regard equations as undirected, they often prove theorems that
require the reverse orientation of a \textit{simp} rule. Higher-order problems
can be tackled, but the success rate is better for first-order problems. Hence,
you may get better results if you first simplify the problem to remove
higher-order features.

Note that problems can be easy for auto and difficult for ATPs, but the reverse
is also true, so don't be discouraged if your first attempts fail. Because the
system refers to all theorems known to Isabelle, it is particularly suitable
when your goal has a short proof from lemmas that you don't know about.

\section{Command Syntax}
\label{command-syntax}

Sledgehammer can be invoked at any point when there is an open goal by entering
the \textbf{sledgehammer} command in the theory file. Its general syntax is as
follows:

\prew
\textbf{sledgehammer} \textit{subcommand\/$^?$ options\/$^?$ facts\_override\/$^?$ num\/$^?$}
\postw

For convenience, Sledgehammer is also available in the ``Commands'' submenu of
the ``Isabelle'' menu in Proof General or by pressing the Emacs key sequence C-c
C-a C-s. This is equivalent to entering the \textbf{sledgehammer} command with
no arguments in the theory text.

In the general syntax, the \textit{subcommand} may be any of the following:

\begin{enum}
\item[$\bullet$] \textbf{\textit{run} (the default):} Runs Sledgehammer on subgoal number
\textit{num} (1 by default), with the given options and facts.

\item[$\bullet$] \textbf{\textit{minimize}:} Attempts to minimize the provided facts
(specified in the \textit{facts\_override} argument) to obtain a simpler proof
involving fewer facts. The options and goal number are as for \textit{run}.

\item[$\bullet$] \textbf{\textit{messages}:} Redisplays recent messages issued by
Sledgehammer. This allows you to examine results that might have been lost due
to Sledgehammer's asynchronous nature. The \textit{num} argument specifies a
limit on the number of messages to display (5 by default).

\item[$\bullet$] \textbf{\textit{available\_atps}:} Prints the list of installed ATPs.
See \S\ref{installation} and \S\ref{mode-of-operation} for more information on
how to install ATPs.

\item[$\bullet$] \textbf{\textit{running\_atps}:} Prints information about currently
running ATPs, including elapsed runtime and remaining time until timeout.

\item[$\bullet$] \textbf{\textit{kill\_atps}:} Terminates all running ATPs.

\item[$\bullet$] \textbf{\textit{refresh\_tptp}:} Refreshes the list of remote
ATPs available at System\-On\-TPTP \cite{sutcliffe-2000}.
\end{enum}

Sledgehammer's behavior can be influenced by various \textit{options}, which can
be specified in brackets after the \textbf{sledgehammer} command. The
\textit{options} are a list of key--value pairs of the form ``[$k_1 = v_1,
\ldots, k_n = v_n$]''. For Boolean options, ``= \textit{true}'' is optional. For
example:

\prew
\textbf{sledgehammer} [\textit{isar\_proof}, \,\textit{timeout} = 120$\,s$]
\postw

Default values can be set using \textbf{sledgehammer\_\allowbreak params}:

\prew
\textbf{sledgehammer\_params} \textit{options}
\postw

The supported options are described in \S\ref{option-reference}.

The \textit{facts\_override} argument lets you alter the set of facts that go
through the relevance filter. It may be of the form ``(\textit{facts})'', where
\textit{facts} is a space-separated list of Isabelle facts (theorems, local
assumptions, etc.), in which case the relevance filter is bypassed and the given
facts are used. It may also be of the form (\textit{add}:\ \textit{facts}$_1$),
(\textit{del}:\ \textit{facts}$_2$), or (\textit{add}:\ \textit{facts}$_1$\
\textit{del}:\ \textit{facts}$_2$), where the relevance filter is instructed to
proceed as usual except that it should consider \textit{facts}$_1$
highly-relevant and \textit{facts}$_2$ fully irrelevant.

\section{Option Reference}
\label{option-reference}

\def\flushitem#1{\item[]\noindent\kern-\leftmargin \textbf{#1}}
\def\qty#1{$\left<\textit{#1}\right>$}
\def\qtybf#1{$\mathbf{\left<\textbf{\textit{#1}}\right>}$}
\def\optrue#1#2{\flushitem{\textit{#1} $\bigl[$= \qtybf{bool}$\bigr]$\quad [\textit{true}]\hfill (neg.: \textit{#2})}\nopagebreak\\[\parskip]}
\def\opfalse#1#2{\flushitem{\textit{#1} $\bigl[$= \qtybf{bool}$\bigr]$\quad [\textit{false}]\hfill (neg.: \textit{#2})}\nopagebreak\\[\parskip]}
\def\opsmart#1#2{\flushitem{\textit{#1} $\bigl[$= \qtybf{bool\_or\_smart}$\bigr]$\quad [\textit{smart}]\hfill (neg.: \textit{#2})}\nopagebreak\\[\parskip]}
\def\opsmartx#1#2{\flushitem{\textit{#1} $\bigl[$= \qtybf{bool\_or\_smart}$\bigr]$\quad [\textit{smart}]\hfill\\\hbox{}\hfill (neg.: \textit{#2})}\nopagebreak\\[\parskip]}
\def\opnodefault#1#2{\flushitem{\textit{#1} = \qtybf{#2}} \nopagebreak\\[\parskip]}
\def\opdefault#1#2#3{\flushitem{\textit{#1} = \qtybf{#2}\quad [\textit{#3}]} \nopagebreak\\[\parskip]}
\def\oparg#1#2#3{\flushitem{\textit{#1} \qtybf{#2} = \qtybf{#3}} \nopagebreak\\[\parskip]}
\def\opargbool#1#2#3{\flushitem{\textit{#1} \qtybf{#2} $\bigl[$= \qtybf{bool}$\bigr]$\hfill (neg.: \textit{#3})}\nopagebreak\\[\parskip]}
\def\opargboolorsmart#1#2#3{\flushitem{\textit{#1} \qtybf{#2} $\bigl[$= \qtybf{bool\_or\_smart}$\bigr]$\hfill (neg.: \textit{#3})}\nopagebreak\\[\parskip]}

Sledgehammer's options are categorized as follows:\ mode of operation
(\S\ref{mode-of-operation}), problem encoding (\S\ref{problem-encoding}), output
format (\S\ref{output-format}), and timeouts (\S\ref{timeouts}).

The descriptions below refer to the following syntactic quantities:

\begin{enum}
\item[$\bullet$] \qtybf{string}: A string.
\item[$\bullet$] \qtybf{bool\/}: \textit{true} or \textit{false}.
\item[$\bullet$] \qtybf{bool\_or\_smart\/}: \textit{true}, \textit{false}, or \textit{smart}.
\item[$\bullet$] \qtybf{int\/}: An integer.
\item[$\bullet$] \qtybf{int\_or\_smart\/}: An integer or \textit{smart}.
\item[$\bullet$] \qtybf{time}: An integer followed by $\textit{min}$ (minutes), $s$ (seconds), or \textit{ms}
(milliseconds), or the keyword \textit{none} ($\infty$ years).
\end{enum}

Default values are indicated in square brackets. Boolean options have a negated
counterpart (e.g., \textit{debug} vs.\ \textit{no\_debug}). When setting
Boolean options, ``= \textit{true}'' may be omitted.

\subsection{Mode of Operation}
\label{mode-of-operation}

\begin{enum}
%\optrue{blocking}{non\_blocking}
%Specifies whether the \textbf{sledgehammer} command should operate synchronously.
%The asynchronous (non-blocking) mode lets the user start proving the putative
%theorem while Sledgehammer looks for a counterexample, but it can also be more
%confusing. For technical reasons, automatic runs currently always block.

\opnodefault{atps}{string}
Specifies the ATPs (automated theorem provers) to use as a space-separated list
(e.g., ``\textit{e}~\textit{spass}''). The following ATPs are supported:

\begin{enum}
\item[$\bullet$] \textbf{\textit{e}:} E is an ATP developed by Stephan Schulz
\cite{schulz-2002}. To use E, set the environment variable
\texttt{E\_HOME} to the directory that contains the \texttt{eproof} executable,
or install the prebuilt E package from Isabelle's download page. See
\S\ref{installation} for details.

\item[$\bullet$] \textbf{\textit{spass}:} SPASS is an ATP developed by Christoph
Weidenbach et al.\ \cite{weidenbach-et-al-2009}. To use SPASS, set the
environment variable \texttt{SPASS\_HOME} to the directory that contains the
\texttt{SPASS} executable, or install the prebuilt SPASS package from Isabelle's
download page. Sledgehammer requires version 3.5 or above. See
\S\ref{installation} for details.

\item[$\bullet$] \textbf{\textit{vampire}:} Vampire is an ATP developed by
Andrei Voronkov and his colleagues \cite{riazanov-voronkov-2002}. To use
Vampire, set the environment variable \texttt{VAMPIRE\_HOME} to the directory
that contains the \texttt{vampire} executable.

\item[$\bullet$] \textbf{\textit{remote\_e}:} The remote version of E runs
on Geoff Sutcliffe's Miami servers \cite{sutcliffe-2000}.

\item[$\bullet$] \textbf{\textit{remote\_vampire}:} The remote version of
Vampire runs on Geoff Sutcliffe's Miami servers. Version 9 is used.

\item[$\bullet$] \textbf{\textit{remote\_sine\_e}:} SInE-E is a metaprover
developed by Kry\v stof Hoder \cite{sine} based on E. The remote version of
SInE runs on Geoff Sutcliffe's Miami servers.

\item[$\bullet$] \textbf{\textit{remote\_snark}:} SNARK is a prover
developed by Stickel et al.\ \cite{snark}. The remote version of
SNARK runs on Geoff Sutcliffe's Miami servers.
\end{enum}

By default, Sledgehammer will run E, SPASS, Vampire, and SInE-E in parallel.
For at most two of E, SPASS, and Vampire, it will use any locally installed
version if available. For historical reasons, the default value of this option
can be overridden using the option ``Sledgehammer: ATPs'' from the ``Isabelle''
menu in Proof General.

It is a good idea to run several ATPs in parallel, although it could slow down
your machine. Running E, SPASS, and Vampire together for 5 seconds yields about
the same success rate as running the most effective of these (Vampire) for 120
seconds \cite{boehme-nipkow-2010}.

\opnodefault{atp}{string}
Alias for \textit{atps}.

\opfalse{overlord}{no\_overlord}
Specifies whether Sledgehammer should put its temporary files in
\texttt{\$ISA\-BELLE\_\allowbreak HOME\_\allowbreak USER}, which is useful for
debugging Sledgehammer but also unsafe if several instances of the tool are run
simultaneously. The files are identified by the prefix \texttt{prob\_}; you may
safely remove them after Sledgehammer has run.

\nopagebreak
{\small See also \textit{debug} (\S\ref{output-format}).}
\end{enum}

\subsection{Problem Encoding}
\label{problem-encoding}

\begin{enum}
\opfalse{explicit\_apply}{implicit\_apply}
Specifies whether function application should be encoded as an explicit
``apply'' operator. If the option is set to \textit{false}, each function will
be directly applied to as many arguments as possible. Enabling this option can
sometimes help discover higher-order proofs that otherwise would not be found.

\opfalse{full\_types}{partial\_types}
Specifies whether full-type information is exported. Enabling this option can
prevent the discovery of type-incorrect proofs, but it also tends to slow down
the ATPs significantly. For historical reasons, the default value of this option
can be overridden using the option ``Sledgehammer: ATPs'' from the ``Isabelle''
menu in Proof General.
\end{enum}

\subsection{Relevance Filter}
\label{relevance-filter}

\begin{enum}
\opdefault{relevance\_thresholds}{int\_pair}{45~95}
Specifies the thresholds above which facts are considered relevant by the
relevance filter. The first threshold is used for the first iteration of the
relevance filter and the second threshold is used for the last iteration (if it
is reached). The effective threshold is quadratically interpolated for the other
iterations. Each threshold ranges from 0 to 100, where 0 means that all theorems
are relevant and 100 only theorems that refer to previously seen constants.

\opdefault{max\_relevant}{int\_or\_smart}{\textit{smart}}
Specifies the maximum number of facts that may be returned by the relevance
filter. If the option is set to \textit{smart}, it is set to a value that was
empirically found to be appropriate for the ATP. A typical value would be 300.

\opsmartx{theory\_relevant}{theory\_irrelevant}
Specifies whether the theory from which a fact comes should be taken into
consideration by the relevance filter. If the option is set to \textit{smart},
it is taken to be \textit{true} for SPASS and \textit{false} for the other ATPs;
empirical results suggest that these are the best settings.

%\opfalse{defs\_relevant}{defs\_irrelevant}
%Specifies whether the definition of constants occurring in the formula to prove
%should be considered particularly relevant. Enabling this option tends to lead
%to larger problems and typically slows down the ATPs.
\end{enum}

\subsection{Output Format}
\label{output-format}

\begin{enum}

\opfalse{verbose}{quiet}
Specifies whether the \textbf{sledgehammer} command should explain what it does.

\opfalse{debug}{no\_debug}
Specifies whether Nitpick should display additional debugging information beyond
what \textit{verbose} already displays. Enabling \textit{debug} also enables
\textit{verbose} behind the scenes.

\nopagebreak
{\small See also \textit{overlord} (\S\ref{mode-of-operation}).}

\opfalse{isar\_proof}{no\_isar\_proof}
Specifies whether Isar proofs should be output in addition to one-liner
\textit{metis} proofs. Isar proof construction is still experimental and often
fails; however, they are usually faster and sometimes more robust than
\textit{metis} proofs.

\opdefault{isar\_shrink\_factor}{int}{1}
Specifies the granularity of the Isar proof. A value of $n$ indicates that each
Isar proof step should correspond to a group of up to $n$ consecutive proof
steps in the ATP proof.

\end{enum}

\subsection{Timeouts}
\label{timeouts}

\begin{enum}
\opdefault{timeout}{time}{$\mathbf{60}$ s}
Specifies the maximum amount of time that the ATPs should spend looking for a
proof. For historical reasons, the default value of this option can be
overridden using the option ``Sledgehammer: Time Limit'' from the ``Isabelle''
menu in Proof General.
\end{enum}

\let\em=\sl
\bibliography{../manual}{}
\bibliographystyle{abbrv}

\end{document}
